Fix the calibrated value of \(\kappa\), fix \(\gamma>d\), and put \(q_0=\min\{\beta,\gamma\}\). We first derive the global smoothness that the estimator will use. At \(t=t_0=c\), \(\mathrm{ass:lower-margin}\) and calibration give
\[
1=c_m c\le P_X\{0<|\tau_P(X)-\theta|\le c\}\le1.
\]
Thus the displayed probability equals one. The density is bounded below by \(2^{-d}\) on the whole cube and \(\tau_P\) is continuous. If \(|\tau_P(x)-\theta|>c\) at some \(x\in\mathcal X\), continuity would give a relatively open neighborhood of positive Lebesgue, and hence positive \(P_X\), measure on which the same strict inequality holds, contradicting the preceding equality. Consequently
\[
|\tau_P(x)-\theta|\le c\quad\text{for every }x\in\mathcal X,
\qquad
P_X\{\tau_P(X)=\theta\}=0,
\qquad
\mathcal D_P(c)=\mathcal X.
\]
The near-band condition is therefore a global \(\gamma\)-Hölder condition.

We next make the Hölder radii explicit. On the fixed cube \([-1,1]^d\), the Taylor-remainder norm has the following elementary nesting property: if \(s\ge q>0\), then
\[
\|g\|_{\mathsf H^q(\mathcal X)}\le C_{\mathrm{emb}}(d,s,q)\|g\|_{\mathsf H^s(\mathcal X)}.
\]
Indeed, derivatives through order \(\lceil q\rceil-1\) are already bounded by the order-\(s\) norm. For pairs at distance at most one, Taylor's formula at the smaller order and the order-\(s\) remainder give the order-\(q\) remainder; for pairs at distance greater than one, bounded derivatives and the fixed diameter \(2\sqrt d\) give the same inequality after enlarging the constant. Applying this first with \((s,q)=(\beta,q_0)\) and then with \((s,q)=(\gamma,q_0)\), and using the calibrated radii \(L_0=L_\tau=1\), yields a fixed constant \(L_*<\infty\) such that
\[
\|\mu_0\|_{\mathsf H^{q_0}(\mathcal X)}\le L_*,
\qquad
\|\tau_P\|_{\mathsf H^{q_0}(\mathcal X)}\le L_*.
\]
Take the continuous version \(\mu_1=\mu_0+\tau_P\), which equals the treated conditional mean \(P_X\)-almost everywhere by \(\mathrm{def:cate}\). Then
\[
\|\mu_1\|_{\mathsf H^{q_0}(\mathcal X)}\le2L_*.
\]
These bounds are uniform over \(P\in\mathcal M_{\mathrm{tr}}(\gamma;\kappa)\). Since \(|Y|\le1\), both continuous regression versions also take values in \([-1,1]\): an excursion outside that interval would, by continuity and the density lower bound, contradict the almost-sure conditional-mean bound.

We now give one estimator that uses only the sample, the known cube, and the frozen constants \(d\), \(q_0\), and \(c_\pi\). Let \(p_0=\lceil q_0\rceil-1\), use \(W_{x,h}\), \(K_{x,h}\), and the shifted-Legendre vector \(r(u)=r_{x,h,p_0}(u)\) from \(\mathrm{def:window-basis}\), and write \(r_x=r(x)\). For \(a\in\{0,1\}\), define
\[
\widetilde Q_a(x)=\frac1n\sum_{i=1}^nK_{x,h}(X_i)r(X_i)r(X_i)^\top\mathbf 1\{A_i=a\},
\qquad
\widetilde R_a(x)=\frac1n\sum_{i=1}^nK_{x,h}(X_i)r(X_i)Y_i\mathbf 1\{A_i=a\}.
\]
Clip the eigenvalues of the symmetric matrix \(\widetilde Q_a(x)\) to
\[
I_a=[2^{-d-1}c_\pi,\,2^{1-d}(1-c_\pi)]
\]
and call the resulting matrix \(\widehat Q_a(x)\). Set
\[
\widehat\mu_a(x)=\operatorname{clip}_{[-1,1]}\!\left\{r_x^\top\widehat Q_a(x)^{-1}\widetilde R_a(x)\right\},
\qquad
\widehat\tau_n^{\mathrm{dir}}(x)=\widehat\mu_1(x)-\widehat\mu_0(x),
\]
\[
\widehat G_n^{\mathrm{dir}}=\{x\in\mathcal X:\widehat\tau_n^{\mathrm{dir}}(x)>\theta\}.
\]
No quantity in these displays uses the unknown \(\pi\), \(\mu_a\), \(f_P\), or any other nuisance belonging to the member law. The maximum and minimum operations defining the windows are continuous in \(x\); the kernel and finite empirical sums are jointly Borel measurable in the sample and \(x\). Spectral clipping is a continuous, hence Borel, map on symmetric matrices, and inversion is continuous on the clipped positive-definite range. It follows that \((Z_1,\ldots,Z_n,x)\mapsto\widehat\tau_n^{\mathrm{dir}}(x)\) is jointly Borel measurable and that every section \(\widehat G_n^{\mathrm{dir}}\) is a Borel set. Thus this is one jointly sample- and \(x\)-measurable rule, common to the entire model class.

For the pointwise analysis, introduce only as proof objects
\[
Q_a(x)=\mathbb E_P\!\left[K_{x,h}(X)r(X)r(X)^\top\mathbf 1\{A=a\}\right],
\qquad e_1(u)=\pi(u),\quad e_0(u)=1-\pi(u).
\]
By \(\mathrm{lem:uniform-design}\), the change of variables \(z=T_{x,h}(u)\), orthonormality of the shifted-Legendre basis, and strict overlap, for every vector \(v\),
\[
2^{-d}c_\pi\|v\|_2^2
\le v^\top Q_a(x)v
\le2^{-d}(1-c_\pi)\|v\|_2^2.
\]
Thus \(Q_a(x)\) lies in the spectral interval \(I_a\), while every \(\widehat Q_a(x)^{-1}\) has operator norm at most \(2^{d+1}c_\pi^{-1}\).

Let \(t_{a,x}\) be the degree-\(p_0\) Taylor polynomial of \(\mu_a\) at \(x\), expressed in the basis \(r\) as \(t_{a,x}(u)=b_{a,x}^\top r(u)\). The uniform \(q_0\)-Hölder radii and \(h/2\le\ell_j\le h\) give
\[
\sup_{u\in W_{x,h}}|\mu_a(u)-t_{a,x}(u)|\le C h^{q_0},
\qquad
r_x^\top b_{a,x}=\mu_a(x).
\]
Moreover,
\[
\sup_{P,x,a}\|b_{a,x}\|_2\le C_b<\infty.
\]
To see the last assertion, write the Taylor polynomial first in the rescaled monomials \(T_{x,h}(u)^\nu\). Each coefficient is a fixed linear combination of derivatives of \(\mu_a\) of order at most \(p_0\), multiplied by products of window lengths and powers of the rescaled location of \(x\). The derivative bounds, \(h\le1\), \(h/2\le\ell_j\le h\), and \(T_{x,h}(x)\in[0,1]^d\) bound all these coefficients. Conversion from the finitely many monomials to the fixed shifted-Legendre basis is an invertible finite-dimensional linear map depending only on \(d\) and \(p_0\), proving the claim. The Taylor remainder also gives
\[
\left\|\mathbb E_P\widetilde R_a(x)-Q_a(x)b_{a,x}\right\|_2\le C h^{q_0}.
\]

All coordinates of \(r\) and \(r_x\) are bounded by a constant depending only on \(d\) and \(p_0\). Since \(K_{x,h}\le Ch^{-d}\), its support has probability at most \(Ch^d\), and \(|Y|\le1\), every centered coordinate \(S\) of either one-observation response summand or Gram summand satisfies
\[
\mathbb E_P S^2\le Ch^{-d},
\qquad
|S|\le Ch^{-d},
\qquad
\mathbb E_P|S|^m\le(Ch^{-d})^{m-2}\mathbb E_P S^2\quad(m\ge2).
\]
Expanding the exponential series under these displayed moment inequalities gives the scalar Bernstein inequality
\[
P\!\left(\left|\frac1n\sum_{i=1}^nS_i\right|>
C\left\{\sqrt{\frac{z}{nh^d}}+\frac{z}{nh^d}\right\}\right)
\le2e^{-z},
\qquad z>0.
\]
Because the polynomial dimension is fixed, a union bound over response coordinates and Gram entries changes only the constants.

It remains to include the empirical-Gram perturbation explicitly. Spectral clipping is the Frobenius projection onto the closed convex set of symmetric matrices with spectrum in \(I_a\). Since \(Q_a(x)\) belongs to that set, the projection variational inequality and Cauchy--Schwarz imply
\[
\|\widehat Q_a(x)-Q_a(x)\|_F
\le\|\widetilde Q_a(x)-Q_a(x)\|_F.
\]
The exact decomposition is
\[
\widehat Q_a^{-1}\widetilde R_a-b_{a,x}
=\widehat Q_a^{-1}\!\left[
\{\widetilde R_a-\mathbb E_P\widetilde R_a\}
+\{\mathbb E_P\widetilde R_a-Q_ab_{a,x}\}
+\{Q_a-\widehat Q_a\}b_{a,x}
\right].
\]
The first term is the response fluctuation, the second is the Taylor bias, and the third is the empirical-Gram perturbation. Combining the inverse bound, the uniform coefficient bound, the Taylor remainder, the projection contraction, and the coordinatewise Bernstein inequalities yields, uniformly in \(P\), \(x\), and \(a\),
\[
P\!\left(|\widehat\mu_a(x)-\mu_a(x)|>
C\left\{h^{q_0}+\sqrt{\frac{z}{nh^d}}+\frac{z}{nh^d}\right\}\right)
\le C_0e^{-c_0z},
\qquad z\ge1.
\]
Final clipping to \([-1,1]\) cannot increase the error because \(\mu_a(x)\in[-1,1]\).

Choose \(h=n^{-1/(2q_0+d)}\) and put
\[
\varepsilon_n=h^{q_0}=n^{-q_0/(2q_0+d)}=(nh^d)^{-1/2}.
\]
For all sufficiently large \(n\), a union bound over the two arms and the inequalities \(\sqrt z\le z\) and \(\varepsilon_n z\le z\), valid for \(z\ge1\), give fixed constants \(B_0,B_1,b_1>0\) such that
\[
P\{|\widehat\tau_n^{\mathrm{dir}}(x)-\tau_P(x)|>B_0z\varepsilon_n\}
\le B_1e^{-b_1z}.
\]
The bound itself holds for every \(z\ge1\); the transfer below uses precisely the admissible range
\[
1\le z\le\frac{c}{2B_0\varepsilon_n}.
\]
Equivalently, after changing constants, for \(2B_0\varepsilon_n\le u\le c\),
\[
P\{|\widehat\tau_n^{\mathrm{dir}}(x)-\tau_P(x)|\ge u\}
\le C\exp\{-c'u/\varepsilon_n\}.
\]

Finally set \(\Delta(x)=|\tau_P(x)-\theta|\). We proved above that \(0\le\Delta(x)\le c\) throughout the cube. A set error at \(x\) implies \(|\widehat\tau_n^{\mathrm{dir}}(x)-\tau_P(x)|\ge\Delta(x)\). Choose a fixed \(A>2B_0\). By the upper margin condition with \(\xi=1\), the contribution from \(0<\Delta(X)\le A\varepsilon_n\) is at most
\[
A\varepsilon_nP_X\{0<\Delta(X)\le A\varepsilon_n\}
\le C\varepsilon_n^2.
\]
For sufficiently large \(n\), choose the unique integer \(J\ge0\) such that
\[
2^JA\varepsilon_n<c\le2^{J+1}A\varepsilon_n.
\]
For \(0\le j<J\), on the shell \(2^jA\varepsilon_n<\Delta(X)\le2^{j+1}A\varepsilon_n\), the pointwise exponential inequality and the upper margin bound give a contribution no larger than
\[
C(2^j\varepsilon_n)^2e^{-c'2^j}.
\]
The final truncated shell \(2^JA\varepsilon_n<\Delta(X)\le c\) has \(2^JA\varepsilon_n\ge c/2\); its contribution is therefore at most \(C e^{-c''/\varepsilon_n}\le C\varepsilon_n^2\). Since \(\sum_{j\ge0}4^je^{-c'2^j}<\infty\), Tonelli's theorem and \(\mathrm{def:weighted-loss}\) now give
\[
\sup_{P\in\mathcal M_{\mathrm{tr}}(\gamma;\kappa)}
\mathbb E_P[d_H(\widehat G_n^{\mathrm{dir}},G_P;P)]
\le C\varepsilon_n^2
=Cn^{-2q_0/(2q_0+d)}.
\]
This is a risk bound for the single data-driven estimator just constructed, not for a law-specific oracle. Taking the infimum over all measurable set estimators in \(\mathrm{def:minimax-risk}\) and substituting \(q_0=\min\{\beta,\gamma\}\) proves both conclusions.